\section{Experimental Setup}
\label{sec:experiments}

\subsection{Training Strategy}

Given the large number of parameters in \arse{} (394.8M), training end-to-end from
scratch would both be computationally prohibitive and likely overfit the moderately
sized combined dataset.  We therefore adopt a two-stage curriculum:

\paragraph{Stage 1 — Frozen-encoder pre-training.}
The \modernbert{} backbone weights are frozen; only the \slpe{} projection, the
cross-attention fusion module, and the \uacc{} classifier are updated (4.9M parameters,
$\approx 1.2\%$ of the total).  This stage teaches the fusion and classification
components to exploit pre-trained semantic representations without distorting them.
The model trained in this stage is saved as \texttt{radar-encoder-freeze} and serves
as the initialisation point for Stage 2.

\paragraph{Stage 2 — End-to-end fine-tuning.}
Starting from the Stage-1 checkpoint, all 399.7M parameters are unfrozen and jointly
optimised.  The lower learning rate (same value as Stage 1) combined with the warm
start prevents catastrophic forgetting of the \modernbert{} pretraining while allowing
the backbone to adapt its representations to the detection task.

\subsection{Hyperparameters}

Table~\ref{tab:hparams} lists the hyperparameters common to all four training
experiments.  We use a linear warm-up for the first 3\% of steps followed by
cosine annealing to zero.  The fused AdamW implementation in PyTorch 2.x
reduces memory overhead and accelerates weight updates.

\begin{table}[t]
  \caption{Hyperparameters shared across all training runs.}
  \label{tab:hparams}
  \centering
  \begin{tabular}{ll}
    \toprule
    Hyperparameter & Value \\
    \midrule
    Optimiser              & AdamW (fused) \\
    Learning rate          & $2 \times 10^{-5}$ \\
    Weight decay           & $1 \times 10^{-2}$ \\
    Training epochs        & 10 \\
    Batch size             & 32 \\
    LR schedule            & Cosine with 3\% linear warm-up \\
    Gradient clipping      & 1.0 \\
    Loss function          & BCEWithLogitsLoss \\
    Positive class weight  & $w^+ = 0.746$ \\
    Max token length       & 512 \\
    Dropout (fusion)       & 0.1 \\
    Dropout (\slpe{})      & 0.1 \\
    Precision              & bf16 (mixed) \\
    Hardware               & NVIDIA A100 (40 GB) \\
    \bottomrule
  \end{tabular}
\end{table}

\subsection{Evaluation Metrics}
\label{subsec:metrics}

Following the PAN-2026 evaluation protocol~\cite{bevendorff:2025}, model performance is
assessed using five complementary metrics:

\begin{description}
  \item[ROC-AUC] Area under the Receiver Operating Characteristic curve; measures
        ranking quality independently of any decision threshold.
  \item[Brier Score (inverted)] $1 - \frac{1}{N}\sum_i (y_i - \hat{y}_i)^2$;
        measures calibration of predicted probabilities.  Higher is better (range
        $[0,1]$).
  \item[$c@1$~\cite{penas:2011}] A variant of accuracy that rewards abstaining on
        uncertain cases by preserving the accuracy of answered questions when
        $|\hat{y}-0.5| < \tau$ is left unanswered.
  \item[$F_1$ macro] Harmonic mean of precision and recall for both classes, averaged.
  \item[$F_{0.5}$ (upright, $F_{0.5u}$)] Weighted $F$ measure with $\beta=0.5$,
        emphasising precision over recall; relevant for applications where false
        positives are more costly than false negatives.
\end{description}

The \emph{mean} over all five metrics is the primary ranking criterion on the
PAN-2026 leaderboard.

\subsection{Training Experiments}

\paragraph{Experiment 1: radar-encoder-freeze.}
Full combined dataset (38,503 train / 18,423 val); encoder frozen; 10 epochs.
This is the baseline frozen-encoder system.

\paragraph{Experiment 2: RADAR (end-to-end).}
Initialised from Experiment 1; all weights unfrozen; same combined dataset; 10 epochs.
This constitutes the full \radar{} system submitted to PAN-2026.

\paragraph{Experiment 3: radar-encoder-freeze-pan26.}
PAN-2026 data only (23,682 train / 3,589 val); encoder frozen; 10 epochs.
Ablation study isolating the contribution of in-domain PAN-2026 data.

\paragraph{Experiment 4: radar-encoder-freeze-raid.}
RAID data only (sampled 30k train / 10k test); encoder frozen; 10 epochs.
Ablation study isolating the contribution of external, adversarially perturbed RAID data.

\subsection{Baseline}

To contextualise neural performance, we implement a classical TF-IDF SVM baseline using
scikit-learn: character and word n-grams ($n \in \{1,2,3\}$, character; $n \in \{1,2\}$,
word), sublinear TF scaling, and a LinearSVC classifier.  The baseline is evaluated on
the same splits as the neural models.
